\documentclass[11pt,a4paper]{article}
\usepackage[utf8]{inputenc}
\usepackage[T1]{fontenc}
\usepackage{amsmath,amssymb,amsfonts,amsthm,mathtools}
\usepackage{graphicx}
\usepackage[margin=1in]{geometry}
\usepackage{fancyhdr}
\usepackage{enumitem}
\usepackage{multicol}
\usepackage{xcolor}
\usepackage{tikz}
\usepackage{subcaption}
\usepackage{kotex}
\usepackage{cancel}
\usepackage[colorlinks=true]{hyperref}
\usepackage{xcolor}
\newenvironment{solution}{%
    \par\noindent\color{blue}\textbf{Answer)}\par
}{%
    \par
}

\pagestyle{fancy}
\fancyhf{}
\renewcommand{\headrulewidth}{0.4pt}
\renewcommand{\footrulewidth}{0.4pt}

\lhead{\textbf{EN5232}}
\chead{\textbf{Climate Dynamics}}
\rhead{\textbf{2026-01}}

\cfoot{-~\thepage~-}

\newcounter{question}

\newcommand{\hwquestion}[2][]{ 
  \stepcounter{question} 
  \vspace{0.6cm} 
  \noindent 
  \ifx\relax#1\relax
  \else
    {\footnotesize(#1)}\\
  \fi
  \noindent \textbf{Q\thequestion. }#2 
  \vspace{0.2cm} 
}

\newtheorem{theorem}{Theorem}[section]

\definecolor{neongreen}{RGB}{57,255,20}
\definecolor{vividmagenta}{RGB}{255,0,255}
\definecolor{brightorange}{RGB}{255,128,0}
\definecolor{vividskublue}{RGB}{0,191,255}
\definecolor{brightgold}{RGB}{255,215,0}
\definecolor{brightlime}{RGB}{191,255,0}
\definecolor{vividpurple}{RGB}{159,0,255}
\definecolor{forestgreen}{RGB}{34,139,34}
\hypersetup{
    colorlinks=true,
    linkcolor=vividskublue,
    citecolor=brightlime,
}

\begin{document}

\providecommand{\hwtitle}{Midterm}
\providecommand{\hwduedate}{2026.04.24}
\providecommand{\studentid}{20251214}
\providecommand{\studentname}{김희수}

\begin{center}
    \Large~\textbf{\hwtitle}
\end{center}

\vspace{0.4cm}

\noindent
\begin{minipage}[t]{0.5\textwidth}
Due: \hwduedate
\end{minipage}
\begin{minipage}[t]{0.5\textwidth}
    \begin{flushright}
    Student ID: \studentid\\
    Name: \studentname
    \end{flushright}
\end{minipage}

\vspace{0.4cm}
\noindent\rule{\textwidth}{0.4pt}



\hwquestion[]{Based on the figures in Chapter 3, compare the amplitudes of the stationary waves in the upper tropospheric circulation in the winter and summer hemispheres.}
\begin{solution}
One of the most notable stationary waves in the upper troposphere is the jet stream in midlatitude. In the winter hemispheres, jet stream strengthens near the right-end of the continents: east Asia and eastern North America. Meanwhile, wind speed of the jet stream decreases. Another key difference between two seasons occurs at Indian Ocean. During JJA, a strong north-east wind blows at Indian ocean, but in DJF, the wind weakens and the direction changes from north-east to south-west. On the pacific ocean, some of jet streams cross the hemisphere and contributes to the southern hemisphere jst stream in DJF, but it disappear in JJA.
\end{solution}

\hwquestion[]{Are any of the features in the climatology of the Earth's skin temperature shown in Figure below explainable without reference to the general circulation?}
\begin{solution}
    Skin temperature largely follows the general circulation; tropic is hot and the poles are cold. However, there are two regions where their temperatures differ from the surroundings: mountain ranges such as the Andes and the Tibet. These places are represented as high altitude where the air temperature is relatively low. 
\end{solution}

\hwquestion[]{Making use of the accompanying schematic of an idealized wave below, explain why the zonal average, $v_E^2$ exhibits a single maximum, midway between equator and pole, flanked by a pair of maxima in $u_E^2$.}
\begin{solution}
Given that idealized wave, the wave displacement is largest at mid-latitude. Thus, the zonal everage of $v_E^2$ exhibits maximum around B. On the other hand, $u_E$ is associated with the meridional derivative of the wave amplitude, so $u_E$ is maximized near the latitude of maximum wave amplitude. Hence the zonal average of $u_E^2$ has two maxima near A and C.
\end{solution}

\end{document}
